\section{Boundary regularity and the proof of the transfer}\label{sec:proof}
\subsection{Finite responses and a compact recurrence class}
\begin{lemma}[A reused row on distinct report coordinates]\label{lem:tensor}
Let $f_{n,k}$ be bounded functions on probability spaces $(Y_k,\sigma_k)$, with a common bound, converging weak-star to $f_k$, for $1\le k\le t$. Then for every $u\in L^1(\bigotimes_{k=1}^t\sigma_k)$,
\[
 \int u\prod_{k=1}^t f_{n,k}(y_k)\,d\sigma
 \longrightarrow \int u\prod_{k=1}^t f_k(y_k)\,d\sigma.
\]
The assertion also holds for nets. In particular, under finite-response domination, every bounded finite-call response is continuous on $\Gamma_M$.
\end{lemma}
\begin{proof}
For a rectangle indicator $u=\prod_k1_{A_k}$ the integral factors into $t$ weak-star pairings. Finite linear combinations of rectangle indicators are dense in $L^1$ of a finite product of probability spaces. The common product bound controls the approximation error uniformly, which proves the assertion, including for nets. Expand a finite-call program into the finitely many label and action words. Each summand is a raw $L^1$ response density paired with the product of the selected report rows; action probabilities multiply it continuously. Continuous bounded readout dependence passes by uniform approximation on the compact decision space. Reuse of a row is harmless because its occurrences are evaluated at distinct report coordinates. It would not be harmless on a diagonal report law not dominated by the product reference.
\end{proof}

\begin{theorem}[An intrinsic criterion for uniform boundary averages]\label{thm:regularity}
For a compact family $\mathcal K$ of stochastic $M\times M$ matrices, the following are equivalent:
\begin{enumerate}
\item $n^{-1}\sum_{k=0}^{n-1}P^k\to\Pi_P$ uniformly on $\mathcal K$;
\item $P\mapsto\Pi_P$ is continuous on $\mathcal K$;
\item $P\mapsto\rank(I-P)$ is locally constant on $\mathcal K$;
\item $\sup_{P\in\mathcal K}h(P)<\infty$.
\end{enumerate}
Moreover, $\{P:h(P)\le H\}$ is compact and
\begin{equation}\label{eq:cesaro}
 \left\|{1\over n}\sum_{k=0}^{n-1}P^k-\Pi_P\right\|_\infty
 \le {2h(P)\over n}.
\end{equation}
This characterizes uniformity for \emph{all} bounded state rewards and all initial labels, not necessary conditions for a single accidentally constant reward.
\end{theorem}
\begin{proof}
The eigenvalue $1$ of a stochastic matrix is semisimple: a nontrivial Jordan block would make its powers unbounded, whereas $\norm{P^n}_\infty=1$. Its other unit-circle eigenvalues are also semisimple, and their Cesaro averages vanish. This proves existence of $\Pi_P$ and the identities
\[
 (I-P)G_P=G_P(I-P)=I-\Pi_P,\quad
 \Pi_PG_P=G_P\Pi_P=0,\quad P\Pi_P=\Pi_PP=\Pi_P.
\]
They also prove uniqueness of $G_P$. Summing the first identity after multiplication by powers of $P$ gives
$\sum_{k=0}^{n-1}(P^k-\Pi_P)=(I-P^n)G_P$, hence \eqref{eq:cesaro}.

Uniform convergence implies continuity of $\Pi_P$ because each finite Cesaro polynomial is continuous. Its trace is the integer dimension of the eigenspace at $1$, so continuity implies local constancy of $\rank(I-P)$. Conversely, near a given matrix of this fixed rank, a small contour around $1$ separates that eigenvalue from all others. Constancy of the nullity and semisimplicity imply that the eigenvalues inside the contour are exactly the fixed multiplicity of $1$. The contour integral of the resolvent therefore makes $\Pi_P$ continuous on this relative neighborhood. Then $(I-P+\Pi_P)^{-1}-\Pi_P$ is continuous as well. A finite subcover of $\mathcal K$ gives a uniform bound on $G_P$, which proves all four implications.

Finally suppose $P_n\to P$ and $\norm{G_{P_n}}_\infty\le H$. Along any convergent subsequence of the bounded matrices $G_{P_n}$, their group-inverse polynomial identities pass to the limit. The limit is the unique group inverse of $I-P$. Thus $h(P)\le H$. The same argument gives continuity of $G$ and $\Pi$ on this closed bounded sublevel. Compactness follows from compactness of all stochastic matrices. These finite-dimensional identities are the classical group-inverse mechanism \cite{Meyer}; the cost-universal criterion records exactly which degenerations the uniform causal theorem excludes.
\end{proof}

\begin{lemma}[Root hitting and the physical time change]\label{lem:roots}
If $h(P)\le H$, one can choose one label in each closed class such that the expected number of boundary transitions to this set, from any label, is at most
$B=(2M+1)H$. For $1\le d_i\le\bar\mu$, the time-changed matrix $S$ in \eqref{eq:timechange} has the same closed classes and class-absorption probabilities as $P$, and
\begin{equation}\label{eq:timegroup}
 h(S)\le4\bar\mu B.
\end{equation}
If $C$ is a closed class and $g_C=\pi_Cr/(\pi_Cd)$, the equation
\begin{equation}\label{eq:poisson}
 (I-P_C)u=r_C-g_Cd_C
\end{equation}
has a solution with $\osc(u)\le2\bar\mu B$.
\end{lemma}
\begin{proof}
Put the transient block first. Since $\Pi_P$ has zero transient columns, the transient block of $G_P$ is $(I-Q)^{-1}$. Its row sums are expected times to the recurrent set and are at most $H$. In a closed class choose $j$ with stationary mass $\pi_C(j)\ge1/M$. The mean hitting-time identity
\[
 \E_i T_j={(G_{P_C})_{jj}-(G_{P_C})_{ij}\over\pi_C(j)}
 \le2MH
\]
follows by applying $(I-P_C)G_{P_C}=I-\Pi_{P_C}$ to its $j$th column and solving the Dirichlet problem with value zero at $j$. The recurrent row block of $G_P$ is $G_{P_C}$, so the bound is inherited. Adding the transient and within-class times proves the root bound.

The matrix $S$ holds at $i$ for a geometric number of steps of mean $d_i$, then makes a $P$ transition. This interpretation includes self-transitions and shows that class absorption is unchanged and that root hitting has expectation at most $\bar\mu B$. On a recurrent class its stationary law is proportional to $\pi_C(i)d_i$, proving \eqref{eq:classratio}.

For completeness a root hitting bound $T_*$ implies $h(P)\le4T_*$. For $\norm{f}_\infty\le1$ set $v=f-\Pi_P f$, so $\norm{v}_\infty\le2$. Stop its accumulated sum on first hitting a chosen root. In a recurrent class $\pi_Cv=0$, so the resulting function, assigned zero at the root, solves $(I-P)u=v$ also at that root. Before the recurrent set it solves the same Dirichlet problem. Its norm is at most $2T_*$. Subtracting $\Pi_Pu$ gives the unique normalized solution $G_Pf$, of norm at most $4T_*$. Apply this to $S$. For \eqref{eq:poisson}, the forcing is centered under $\pi_C$ and has norm at most $\bar\mu$; the same stopped sum, now inside $C$, has oscillation at most $2\bar\mu B$.
\end{proof}

\subsection{From boundary cycles to actual acquisitions}
\begin{proof}[Proof of Theorem~\ref{thm:main}]
By Lemma~\ref{lem:tensor}, finite-call truncated cycle rewards, lengths and exit probabilities are continuous. Truncation changes expected reward and length by at most $a(T)$. It changes exit probabilities by at most $\PP(\tau>T)\le a(T)/(1)$ for integer $\tau>T$, using a truncation one step earlier if necessary. Thus $P,d,r$ are continuous. The closed-sublevel conclusion of Theorem~\ref{thm:regularity} makes \eqref{eq:regularclass} compact. Lemma~\ref{lem:roots} bounds $G_S$ uniformly, so $\Pi_S$ and $g_\theta$ are continuous on this class.

For a fixed model and program, the finite boundary chain enters a closed class almost surely. Successive returns to a chosen label in that class are regenerative for the \emph{combined} physical experiment and its retained label, with integrable duration and reward. The usual renewal-reward argument, or its stopped-sum proof, gives the almost-sure limit $\pi_C r/\pi_Cd$. The expectation is \eqref{eq:classratio}, since all averages are bounded by one. This use of a return label in a proof does not physically erase or constrain the observer at a cycle boundary. The renewal facts used here are standard \cite{Asmussen}.

We prove the quantitative comparison without a stationarity assumption. First suppose the initial boundary label is in a closed class $C$. At its cycle boundaries, the increments
\[
 R_k-g_C\tau_k+u(I_{k+1})-u(I_k)
\]
have conditional mean zero by \eqref{eq:poisson}. Let $K$ be the first cycle index whose cumulative length exceeds $N$. Then $K\le N+1$. The event that cycle $k$ is reached is determined before its innovations, so the expected stopped sum of these centered increments is zero, by a finite sum of conditional expectations. Removing the portion after call $N$ costs at most the excess length, because the centered per-call loss lies in $[-1,1]$. The bias costs at most $\osc(u)$.

Here is the overshoot estimate, with no iid assumption. If $s$ is a renewal epoch not exceeding $N$, the last cycle contributes $(\tau-(N-s))_+$; condition on the entering label and the entire past at $s$. At most one boundary occurs at each integer $s$, since durations are positive integers. Therefore
\begin{equation}\label{eq:overshoot}
 \E(S_K-N)
 \le\sum_{s=0}^N\PP(s\text{ is a boundary})\,
       \sup_i\E_i(\tau-(N-s))_+
 \le\sum_{t=0}^N a(t).
\end{equation}
This proves the bound with $2\bar\mu B$ in place of $3\bar\mu B$ when starting recurrently.

From a transient label, wait in the proof until the first chosen closed-class root. Its expected physical time is at most $\bar\mu B$. The discrepancy between its prefix reward and the eventual class gain times that prefix, truncated at $N$, is at most that time. Conditional on the reached root and its arrival time not exceeding $N$, the preceding recurrent estimate applies to the remaining horizon; its overshoot bound is at most the displayed full-horizon bound. The eventual class has probabilities $w_C$. Adding the prefix gives \eqref{eq:bN}. In particular $b_N\to0$, because Cesaro averages of $a(t)\to0$ vanish.

The Abel comparison uses only bounded per-call losses. The identity
$\E\ell_N=NJ_N-(N-1)J_{N-1}$ gives, after multiplying by $(1-\beta)\beta^{N-1}$ and summing,
\begin{equation}\label{eq:abel}
 J_\beta=(1-\beta)^2\sum_{N\ge1}N\beta^{N-1}J_N.
\end{equation}
These coefficients are nonnegative and sum to one. Hence \eqref{eq:bridges}; the right side tends to zero by splitting off a fixed finite range. Under the moment assumption, fix a model and program and denote its conditional lifetime from label $i$ by $\tau_i$. Instead of interchanging a supremum with a sum, use the finite number of entering labels in \eqref{eq:overshoot}:
\[
 \sum_{s=0}^N\max_i\E_i(\tau_i-s)_+
 \le\sum_{i=1}^M\E_i[\tau_i\min\{\tau_i,N+1\}]
 \le MK(N+1)^{1-p}\le2^{1-p}MKN^{1-p}.
\]
The factor $M$ can be removed if one common stochastic dominating lifetime has this moment bound. A conditional moment bound alone does not justify that removal.
For a geometric variable $G$ on $\N$ with success probability $1-\beta$, concavity gives
$(1-\beta)\E G^{1-p}\le(1-\beta)^p$. Applying this to \eqref{eq:abel} proves \eqref{eq:moment-rate}. Taking supremum over models and infimum over the \emph{same} class preserves each uniform bound.

The worst-model average is lower semicontinuous on a compact class, so it attains its minimum. For a proposed level $v$, each set
\[
 \{\gam\in\Gamma_M:h(P_\theta^\gam)\le H,
                         \ g_\theta(\gam)\le v\}
\]
is closed in the compact ambient program space: any convergent net in it has bounded group inverse, continuous time-change projection and the same inequalities at the limit. Finite intersection now proves \eqref{eq:attain}, including the recurrence constraints in the finite subfamilies. There is no exchange of infimum and supremum. Uniform comparisons and lower semicontinuity give the cluster-point assertion.

Finally, for finitely many models each particular program has a finite maximum of the finitely many numbers $h(P_\theta)$. The union of the integer sublevels is therefore all of $\Gamma_M$. Taking infima proves \eqref{eq:exhaustion}. If an unrestricted optimizer exists it lies in some level; conversely a level attaining the unrestricted value has an optimizer by part (i). This proves the attainment equivalence without a claim of uniformity as $H\to\infty$. Part (v) is proved in Sections~\ref{sec:transfer}--\ref{sec:sharp}: those arguments use only parts (i)--(iii) for existence and time comparison, establish their own geometric and defect inequalities, and then combine them with \eqref{eq:bridges}. There is no use of part (v) in their proofs.
\end{proof}
